\section{局限与未来工作}\label{sec:limitations}

形式检查之外，当前模型在特殊文字技巧和典故事件对应上仍有明显不足。以下保留最终权重的代表性案例，并据此提出后续方向。

\subsection{典型能力局限}

\textbf{字形线索尚未形成特殊对应。}对“山河湖海波涛涌”和“江河湖海波涛涌”，模型在有无规则时均输出“日月星辰气象新”。其中“山”不含氵，而“江河湖海波涛涌”七字均含氵；输出没有回应后一上联的统一偏旁构造。ChineseBERT的字形信息进入了词语表示，但现有监督与解码并未要求生成具有对应构字关系的整句。使用字形特征因此不等于掌握偏旁、拆字或谐音技巧。

\Needspace{9\baselineskip}
\textbf{相关词语尚未构成典故关系。}“商君徙木重诚信”的无规则下联为“商贾兴邦重信诚”，规则约束后为“仁者爱民多善良”。“徙木”出自商鞅兑现赏金以取信于民的事件\cite{shiji}；前者混淆“商君”和“商贾”，后者转为一般道德表述，均未构造相称的历史事件。类似地，常见意象之间的局部联系也不能保证整联内容和意境相对。规则补强能够减少形式违规，却无法提供这些关系所需的文化知识。

本报告的结构通过率不覆盖完整传统联律，\href{https://github.com/mozillazg/python-pinyin}{现代读音代理}也未解决古入声、语境多音及节奏问题。初始与增强配置同时改变了数据规模和模型容量，整体差异不单独归因于架构；消融属于固定权重推理屏蔽。反向评分及规则开关的逐句chrF差区间包含零，统计结果未覆盖重新训练、多种子和专家文学评价的不确定性。赏析分支还受到AI标签与模板化表达的影响。

\subsection{后续研究方向}

典故增强是后续工作的重点。已调研的ACPG研究将历史典故与诗歌建立关联，并通过典故选择模块及分阶段学习支持知识条件生成\cite{acpg}。可据此整理带出处、人物事件和寓意的典故表示，先确定与上联相称的事件，再将其作为生成条件；评价分别检查典故事实、事件关系与表达自然度。该方向尚未纳入本报告的最终模型，未来应先验证正确典故条件能否被生成器有效使用，再比较自动选择与人工给定条件。

偏旁、拆字和音韵技巧可通过专门标注与分层测试进一步研究，将字形组织和语境读音转化为明确的训练目标。评价则增加词组对应及意境关系的人工作品分析，分别衡量模型生成和规则补强的贡献，使优化目标与对联的文学要求相衔接。
